\section{Electric transport and the complete vacuum static base}
\label{stat:section}

The stationary development supplied by the preceding section need not
initially have a timelike Killing field everywhere.  We first transport
the electric alignment from the boundary, and then prove staticity
without making any regularity assumption on the zero set of the Killing
norm.  The charged case is reduced to vacuum by an explicit change of
potential.  In the zero-charge case, a complete Riemannian circle
extension excludes an additional interior zero set.  No black-hole
uniqueness theorem is used in this section.

Throughout this section, $k=n-1$, $c_n^2=k(k-1)/2$, and all spatial
norms, divergences, and volume forms without a subscript refer to the
original metric $g$.  Write $\Omega^\circ=\Omega\setminus S$.  We use
Proposition~\ref{var:stationary-data}, whose hypotheses are the
connected nonextremal equality hypotheses of the main theorem.  In
particular, it supplies smooth functions and fields $u,X,p$ on the
original exterior, with
\begin{equation}
 u>0\quad\hbox{on }\Omega^\circ,\qquad
 u\geq |X|,\qquad
 dp=(k-1)uE^\flat,\qquad
 d(i_X\alpha)=0,\qquad \alpha=i_E\vol_g.
 \label{stat:interface-fields}
\end{equation}
The potential tends to zero at infinity.  On
$\mathbb R_z\times\Omega^\circ$ the stationary Lorentz metric and
Faraday form are
\begin{equation}
 \mathbf g=-u^2dz^2+g_{ij}(dx^i+X^i dz)(dx^j+X^j dz),\qquad
 F=c_n\,dz\wedge uE^\flat.
 \label{stat:stationary-metric}
\end{equation}
The Killing field is $\xi=\partial_z$, and its future unit normal to a
constant-$z$ slice is $\mathbf n=(\partial_z-X)/u$.  Both Maxwell
equations hold.  If $\mathbf G$ is the Einstein tensor and $\mathsf S$
the electromagnetic stress tensor with the normalization fixed in the
problem, the remaining matter tensor satisfies
\begin{equation}
 \mathbf G-\mathsf S
   =\mu_m u^{-2}\xi^\flat\otimes\xi^\flat,
 \qquad \mu_m=0\quad\hbox{where }u>|X|.
 \label{stat:remaining-matter}
\end{equation}

Set
\begin{equation}
 W=u^2-|X|^2=-\mathbf g(\xi,\xi)\geq0,
 \qquad
 \kappa=\frac{k-1}{2r_h}
       \left(1-\frac{Q^2}{r_h^{2(k-1)}}\right)>0.
 \label{stat:W-kappa}
\end{equation}
The boundary conclusions of Proposition~\ref{var:stationary-data} are
\begin{equation}
 X=u\nu,\qquad
 \partial_\nu u+K(X,\nu)=\kappa\quad\hbox{on }S,
 \qquad W|_S=0.
 \label{stat:boundary-laws}
\end{equation}
Either $u>0$ on all of $S$, in which case
$dW=2\kappa X^\flat$ there, or $u=0$ on all of $S$.  In the latter
case, in the original distance $s$ into the exterior from $S$,
\begin{equation}
 u=s a,\qquad X=s^2Y,\qquad a|_S=\kappa,
 \label{stat:zero-lapse-factors}
\end{equation}
with smooth $a,Y$.  Thus $W>0$ on a deleted boundary collar in either
case.  On the end, for a fixed
\begin{equation}
 \frac{n-2}{2}<q_0<\min(q,n-2),
 \label{stat:end-exponent}
\end{equation}
we have
\begin{equation}
 (u,X)=(b^0,b)+O_2(r^{-q_0}),\qquad
 (b^0)^2-|b|_\delta^2=1,\qquad
 p=O_2(r^{2-n}).
 \label{stat:end-interface}
\end{equation}
Decreasing $q_0$ while retaining~\eqref{stat:end-exponent} is harmless.
These are the only conclusions from the preceding section used below.

\subsection{Boundary alignment and electric transport}

\begin{lemma}[Global electric alignment]
\label{stat:electric-alignment}
The potential $p$ is constant on $S$, and
\begin{equation}
 X\wedge E=0\quad\hbox{throughout }\Omega.
 \label{stat:alignment}
\end{equation}
\end{lemma}

\begin{proof}
Suppose first that $u>0$ on $S$.  The stationary metric is smooth and
nondegenerate up to the hypersurface $\mathbb R\times S$.  Equations
\eqref{stat:boundary-laws} imply that $\xi$ is a null normal to this
hypersurface: it is null, is tangent to it, and is orthogonal to all its
tangent vectors.  Its null second fundamental form vanishes, since its
flow preserves the metric induced on the transverse sections.  The
null expansion equation consequently gives
$\operatorname{Ric}_{\mathbf g}(\xi,\xi)=0$.  This assertion only uses
the one-sided smooth jets on the hypersurface.  The nonaffine
normalization of its null generators contributes a multiple of the
expansion, which is zero.

The extra tensor in~\eqref{stat:remaining-matter} has zero contraction
with $(\xi,\xi)$ there.  Maxwell null positivity therefore gives
\[
 0=2\lvert i_\xi F\rvert_{\mathbf g}^2.
\]
A covector annihilating a null vector has nonnegative Lorentz norm,
and its norm is zero precisely when it is proportional to that null
vector's metric dual.  Thus $i_\xi F$ annihilates the tangent space of
the horizon.  Since $i_\xi F=c_nuE^\flat$, the tangential component of
$E$ on $S$ is zero.  Equations~\eqref{stat:interface-fields} and
\eqref{stat:boundary-laws} give $X\wedge E=0$ and $dp|_{TS}=0$.
If $u=0$ on $S$, these two assertions follow immediately from $X=0$
and $dp=(k-1)uE^\flat$.  Connectedness of $S$ gives a single boundary
value, denoted by $p_b$.

The divergence-free constraint and Cartan's formula give
\[
 0=d(i_X\alpha)=\mathcal L_X(i_E\vol_g)
   =i_{[X,E]+(\operatorname{div}X)E}\vol_g.
\]
In particular,
\begin{equation}
 [X,E]+(\operatorname{div}X)E=0,
 \qquad \mathcal L_E(X\wedge E)=0.
 \label{stat:electric-bracket}
\end{equation}
Alignment therefore propagates along every nonzero electric
trajectory once it is known at one point of that trajectory.

Here is the measure argument ensuring sufficiently many boundary
hits.  Let $\Pi$ be a relatively compact regular level patch
$p=p_0\ne0$ in the interior.  Follow $V=\operatorname{sgn}(p_0)E$
from this patch until its first boundary hit, with maximal time
$T(y)\in(0,\infty]$.  Along each trajectory $|p|$ increases strictly,
because
\[
 dp(E)=(k-1)u|E|^2>0
\]
where $E\ne0$.  All these trajectories remain in a fixed compact
subset of the original exterior: outside a sufficiently large
coordinate sphere, $|p|<|p_0|$.  The trajectory map
$(t,y)\mapsto\Phi_t(y)$ is injective for $0<t<T(y)$, by uniqueness of
the flow and strict monotonicity of $p$.

Let $d\mathfrak f$ be the positive transverse flux measure
$i_V\vol_g$ on $\Pi$, oriented appropriately.  Since
$\operatorname{div}V=0$, the volume element in this flow tube is
$dt\,d\mathfrak f$.  Consequently
\begin{equation}
 \int_\Pi T(y)\,d\mathfrak f(y)
 \leq \operatorname{Vol}_g(K_\Pi)<\infty,
 \label{stat:flow-time-volume}
\end{equation}
where $K_\Pi$ is the fixed compact region just described.  Hence
$T(y)<\infty$ for almost every $y$ in flux measure.  A finite-time
trajectory cannot terminate in the interior of that compact region.
This follows by smoothly extending the vector field locally and using
ordinary differential equation existence.  Nor can it reach a zero
of $E$ at a finite time, by uniqueness.  It therefore hits $S$ with
$E\ne0$.

Equation~\eqref{stat:electric-bracket} transports the boundary
alignment back to almost every point of $\Pi$.  The flux density is
smooth and strictly positive on this regular patch, so continuity
gives alignment on the whole patch.  Varying the patch proves it at
every point with $p\ne0$ and $E\ne0$.  If $p=0$ and $E\ne0$ at an
interior point, the derivative of $p$ along $E$ is nonzero there,
so such a point is a limit of the preceding points.  At zeros of $E$
alignment is automatic.  This proves~\eqref{stat:alignment}.
\end{proof}

\subsection{Killing identities and a twist estimate through zero sets}

On the open set $\mathcal U=\{W>0\}$ define
\begin{equation}
\begin{gathered}
 \lambda=\sqrt W,\qquad
 h_s=g+W^{-1}X^\flat\otimes X^\flat,\\
 C=h_s^{-1}=g^{-1}-u^{-2}X\otimes X,\qquad
 A_s=W^{-1}X^\flat,\qquad F_s=dA_s.
\end{gathered}
 \label{stat:orbit-quantities}
\end{equation}
Here $C$ is written as a contravariant tensor on the original slice;
it extends smoothly and is positive semidefinite everywhere in the
interior.  The stationary metric is
\begin{equation}
 \mathbf g=-\lambda^2(dz-A_s)^2+h_s
 \quad\hbox{on }\mathbb R\times\mathcal U.
 \label{stat:orbit-splitting}
\end{equation}

\begin{lemma}[Killing curvature and flux identities]
\label{stat:killing-identities}
Put $D_{ab}=\boldsymbol\nabla_a\xi_b$, where
$\boldsymbol\nabla$ is the Lorentz connection.  Then
\begin{equation}
 \operatorname{Ric}^{a}{}_{b}\xi^b
      =-(k-1)^2|E|^2\xi^a,
 \label{stat:ricci-eigenvalue}
\end{equation}
and, throughout the original open exterior,
\begin{equation}
 u^{-1}\operatorname{div}(uC\,dW)
   =-2D_{ab}D^{ab}+2\operatorname{Ric}(\xi,\xi).
 \label{stat:killing-wave}
\end{equation}
The smooth antisymmetric tensor
$\widetilde Q^{ij}=u(d\xi^\flat)^{ij}$ satisfies
$\nabla_i\widetilde Q^{ij}=0$.  On $\mathcal U$ it is
\begin{equation}
 Q^{ij}=uW C^{ia}C^{jb}(F_s)_{ab}.
 \label{stat:twist-current}
\end{equation}
\end{lemma}

\begin{proof}
In spacetime dimension $k+2$, the electromagnetic trace reversal is
\[
 \operatorname{Ric}_{ab}
   =2F_{ac}F_b{}^c-\frac1k(F_{cd}F^{cd})\mathbf g_{ab}
\]
apart from the matter term in~\eqref{stat:remaining-matter}, which is
trace-free wherever it is nonzero.  In an orthonormal frame whose
timelike vector is $\mathbf n$, the electric magnitude is $c_n|E|$.
The mixed Ricci eigenvalue on the plane spanned by $\mathbf n$ and
$E$ is therefore
\[
 -2c_n^2\left(1-\frac1k\right)|E|^2
   =-(k-1)^2|E|^2.
\]
Lemma~\ref{stat:electric-alignment} places
$\xi=u\mathbf n+X$ in this plane wherever $E\ne0$.  The additional
matter tensor annihilates $\xi$, since it is supported where
$\mathbf g(\xi,\xi)=0$.  The same formula holds at $E=0$.
This proves~\eqref{stat:ricci-eigenvalue}, including
\begin{equation}
 \operatorname{Ric}(\xi,\xi)=(k-1)^2|E|^2W.
 \label{stat:ricci-null-contraction}
\end{equation}

For a Killing field, $D$ is skew and
$\boldsymbol\nabla_a(d\xi^\flat)^{ab}
=-2\operatorname{Ric}^{b}{}_{a}\xi^a$.
Applying the Killing wave identity to $W=-\mathbf g(\xi,\xi)$
gives~\eqref{stat:killing-wave}: a stationary scalar has wave
operator $u^{-1}\operatorname{div}(uC\,d\,\cdot)$ because the
spatial inverse block of $\mathbf g$ is $C$ and
$\vol_{\mathbf g}=u\,dz\wedge\vol_g$.
The spatial components of the right side of the Killing divergence
identity vanish by~\eqref{stat:ricci-eigenvalue}.  Antisymmetry
removes the remaining Christoffel term, and hence
$\nabla_i\widetilde Q^{ij}=0$.
Finally,
\[
 \xi^\flat=-W(dz-A_s),\qquad
 d\xi^\flat=-dW\wedge(dz-A_s)+WF_s.
\]
Raising both spatial indices eliminates the first summand and gives
\eqref{stat:twist-current}.
\end{proof}

\begin{lemma}[Vanishing twist]
\label{stat:vanishing-twist}
The two-form $F_s$ vanishes on every component of $\mathcal U$.
No regularity of the interior set $\{W=0\}$ is required.
\end{lemma}

\begin{proof}
The interior zero set
$Z_0=\{W=0\}\cap\Omega^\circ$ is compact, since $W\to1$ at
infinity and is positive on a deleted boundary collar.  On a fixed
relatively compact neighborhood of $Z_0$, both $u$ and $|X|$ have
positive lower bounds.  Put $a=|X|$, $e=X/a$, and $T=e^\perp$ there.
The eigenvalues of $C$ are one on $T$ and $W/u^2$ on $e$.
Smooth nonnegativity of $W$ implies the local estimate
\begin{equation}
 |dW|^2\leq C_0W.
 \label{stat:glaeser}
\end{equation}
For completeness, this follows by restricting to short coordinate
line segments, using a bound for the Hessian, and minimizing the
resulting quadratic Taylor upper bound at a point where the function
is nonnegative.  A finite cover gives a uniform constant on the
chosen neighborhood.

Use a slice-orthonormal frame $(e,e_A)$ and complete it by the future
normal $\mathbf n$.  Since $\xi=u\mathbf n+ae$,
\[
 D_{i\mathbf n}=-\frac{W_i/2+aD_{ie}}u.
\]
An explicit expansion of the Lorentz norm gives
\begin{align}
 -2D_{ab}D^{ab}
 ={}&\frac{|dW|^2}{u^2}
     +\frac{4a}{u^2}\sum_A W_A D_{Ae}
     -\frac{4W}{u^2}\sum_A D_{Ae}^2
     -4\sum_{A<B}D_{AB}^2.
 \label{stat:D-expansion}
\end{align}
The original fields and their derivatives are bounded on this fixed
neighborhood.  Equations~\eqref{stat:glaeser},
\eqref{stat:ricci-null-contraction}, and
\eqref{stat:killing-wave} therefore imply
\begin{equation}
 u^{-1}\operatorname{div}(uC\,dW)
       \leq C_1\bigl(W+|d_TW|\bigr).
 \label{stat:W-inequality}
\end{equation}

Choose a smooth compactly supported $\eta$ in this neighborhood.
For $\delta>0$, test~\eqref{stat:W-inequality} with
$\eta^2/(W+\delta)$.  There is no singular integration here: $C,W$
are smooth and the denominator is positive.  Integration by parts
gives, with $L W=u^{-1}\operatorname{div}(uC\,dW)$,
\begin{align*}
 I_\delta
 &:=\int u\eta^2\frac{|dW|_C^2}{(W+\delta)^2}
  =\int\frac{u\eta^2LW}{W+\delta}
     +2\int\frac{u\eta\,C(d\eta,dW)}{W+\delta},\\
 |dW|_C^2&=|d_TW|^2+\frac W{u^2}\bigl(dW(e)\bigr)^2.
\end{align*}
Since $0\leq C\leq g^{-1}$, Young's inequality absorbs both the
cutoff term and the $|d_TW|/(W+\delta)$ term.  Explicitly, each is
at most $I_\delta/4$ plus a constant times
$\int u(\eta^2+|d\eta|^2)$.  The remaining
$W/(W+\delta)$ term is bounded.  Consequently
\begin{equation}
 I_\delta\leq C_\eta\quad(0<\delta<1),\qquad
 \int_{\mathcal U}\eta^2|d_T\log W|^2<\infty.
 \label{stat:log-energy}
\end{equation}
Take finitely many such cutoffs equal to one near $Z_0$.

We record the exact flux used at the zero-set cutoff.  Contract
\eqref{stat:twist-current} with $A_s$, and lower the remaining index
using $g$.  Since $C X=(W/u^2)X$, direct substitution of
$F_s=d(W^{-1}X^\flat)$ gives
\begin{equation}
 J^i=Q^{ij}(A_s)_j,\qquad
 J^\flat_i=\frac1u\left(
        X^j(dX^\flat)_{ij}-|X|^2(d_T\log W)_i\right).
 \label{stat:exact-twist-flux}
\end{equation}
In particular $J$ is exactly orthogonal to $X$, and near $Z_0$
\begin{equation}
 |J|\leq C\bigl(1+|d_T\log W|\bigr).
 \label{stat:twist-flux-bound}
\end{equation}

Choose a smooth cutoff $\chi:[0,\infty)\to[0,1]$, zero on $[0,1]$
and one on $[2,\infty)$.  A zero-set cutoff $z_\delta$ can be chosen
to equal $\chi(W/\delta)$ near $Z_0$ and one outside the fixed
neighborhood.  Indeed place the spatial localization transition on
a compact set disjoint from $Z_0$; $W$ has a positive lower bound
there, so for all sufficiently small $\delta$ no transition
derivative remains.  On the support of $dz_\delta$, transversality
in~\eqref{stat:exact-twist-flux} yields
\begin{equation}
 |J(dz_\delta)|
 \leq C\frac W\delta
       \left(|d_T\log W|+|d_T\log W|^2\right)
 \leq C\left(1+|d_T\log W|^2\right).
 \label{stat:zero-cutoff-error}
\end{equation}
This is supported where $\delta<W<2\delta$.  Its integral tends to
zero by~\eqref{stat:log-energy} and dominated convergence.  In
particular, an open portion of $Z_0$ causes no omitted contribution:
all test fields used below vanish on a neighborhood of it.

Next choose a distance cutoff $\beta_\varepsilon$ at $S$, zero for
$s\leq\varepsilon$ and one for $s\geq2\varepsilon$, with derivative
bounded by $C/\varepsilon$.  If $u>0$ on $S$, the boundary laws give
\[
 W\asymp s,\qquad e=\nu_s+O(s),\qquad d_TW=O(s),
\]
where $\nu_s=\nabla s$.  Thus $J=O(1)$ and
$J\cdot\nu_s=O(s)$ by its exact transversality.  The resulting
boundary error is $O(\varepsilon)$.  If $u=0$ on $S$, write
\[
 W=s^2w,\qquad w=a^2-s^2|Y|^2>0
\]
using~\eqref{stat:zero-lapse-factors}.  Then
$A_s=Y^\flat/w$, $F_s$, and $C$ are smooth up to $S$.
Equation~\eqref{stat:twist-current} gives $Q=O(s^3)$ and
$J=O(s^3)$, so the boundary error is $O(\varepsilon^3)$.

It remains to make the test field decay at infinity.  In the end
coordinates let $A_\infty$ be the constant limit of $A_s$.  Choose
a smooth function
\[
 \psi(x)=\rho(r)(A_\infty)_ix^i,
\]
where $\rho$ is zero near the compact core and one far out.  Its
support is chosen disjoint from the boundary and zero-set cutoff
regions.  Set $B=A_s-d\psi$ on $\mathcal U$.  Then $dB=F_s$ and
\begin{equation}
 B=O_1(r^{-q_0}),\qquad Q=O(r^{-1-q_0}).
 \label{stat:twist-end-decay}
\end{equation}
There is no bulk term from the transition of $\psi$, since
$Q^{ij}\nabla_i\nabla_j\psi=0$ by antisymmetry.  If $\theta_R$ is
one for $r\leq R$, zero for $r\geq2R$, and
$|d\theta_R|\leq C/R$, its error is bounded by
\begin{equation}
 \left|\int Q^{ij}B_j\partial_i\theta_R\right|
       \leq C R^{n-2-2q_0}\longrightarrow0.
 \label{stat:twist-infinity-error}
\end{equation}

Test $\nabla_i\widetilde Q^{ij}=0$ with
$z_\delta\beta_\varepsilon\theta_R B_j$, extended by zero where
necessary.  This is a smooth compactly supported test in the
interior.  Antisymmetry gives
\begin{align}
 \int z_\delta\beta_\varepsilon\theta_R\,\mathcal T_{\mathrm{tw}}
   &=-\int Q^{ij}B_j
         \partial_i(z_\delta\beta_\varepsilon\theta_R),
 \label{stat:twist-test}\\
 \mathcal T_{\mathrm{tw}}
   &:=\frac12 Q^{ij}(F_s)_{ij}
    =\frac{uW}{2}C^{ia}C^{jb}(F_s)_{ij}(F_s)_{ab}\geq0.
 \label{stat:twist-density}
\end{align}
The three error estimates are independent of the other cutoffs,
which lie between zero and one.  Take $R\to\infty$,
$\varepsilon\downarrow0$, and $\delta\downarrow0$, or use a
diagonal sequence.  Fatou's lemma in~\eqref{stat:twist-test} gives
$\int_{\mathcal U}\mathcal T_{\mathrm{tw}}=0$; its integrability need not be
assumed beforehand.  Since $C$ is positive definite on
$\mathcal U$, $F_s=0$ there.
\end{proof}

\subsection{The charged polynomial and the regular base attachment}

\begin{lemma}[A global polynomial when the charge is nonzero]
\label{stat:charged-polynomial}
If $Q\ne0$, then $p_b\ne0$, $0<|p_b|<1$, and
\begin{equation}
 W=1+p^2+d_0p
   =\left(1-\frac p{p_b}\right)(1-p_bp),
 \qquad d_0=-\frac{1+p_b^2}{p_b}.
 \label{stat:W-polynomial}
\end{equation}
Moreover $p$ lies strictly between $0$ and $p_b$ in the interior,
and $W>0$ everywhere in $\Omega^\circ$.
\end{lemma}

\begin{proof}
The Maxwell equation in~\eqref{stat:interface-fields} gives
\begin{equation}
 \operatorname{div}(u^{-1}\nabla p)=0.
 \label{stat:potential-equation}
\end{equation}
Its coefficients are smooth and elliptic on every compact subset
of the interior.  A positive maximum above both boundary and end
values, or a negative minimum below them, would be attained in such
a subset.  The strong maximum principle therefore applies without
requiring a uniform boundary bound for $u^{-1}$.  If $p_b=0$, it
would give $p=0$ and then $E=0$, contrary to $Q\ne0$.  Otherwise it
gives the stated strict interval for $p$.

There is a useful flux which is smooth even at interior zeros of
$W$:
\begin{equation}
 \mathcal D^i=u(d\xi^\flat)^{iz}
  =\frac1u\left(\nabla^iW+X^j(dX^\flat)^i{}_j\right).
 \label{stat:global-Killing-flux}
\end{equation}
The Killing divergence identity and
\eqref{stat:ricci-eigenvalue} imply
\begin{equation}
 \operatorname{div}\mathcal D=2u^{-1}|dp|^2,
 \qquad \mathcal D\cdot X=u^{-1}dW(X).
 \label{stat:global-flux-laws}
\end{equation}
On $\mathcal U$, Lemma~\ref{stat:vanishing-twist} and
\eqref{stat:global-Killing-flux} also give
\begin{equation}
 \mathcal D=\frac uW C\,dW.
 \label{stat:positive-Killing-flux}
\end{equation}

Choose $d_0$ as in~\eqref{stat:W-polynomial}, and put
\[
 P=W-1-p^2-d_0p,\qquad
 j=\mathcal D-(2p+d_0)u^{-1}\nabla p.
\]
Equations~\eqref{stat:potential-equation} and
\eqref{stat:global-flux-laws} give $\operatorname{div}j=0$.
On $\mathcal U$, alignment implies
$C\,dp=(W/u^2)\nabla p$; hence
\begin{equation}
 j=\frac uW C\,dP,
 \qquad dP(j)\geq u^{-1}|dP|^2.
 \label{stat:polynomial-energy}
\end{equation}
Indeed the axial eigenvalue of $(u/W)C$ is $1/u$, and its
transverse eigenvalues are $u/W\geq1/u$.  This also covers points
where $X=0$.  At an interior zero of $W$, $dW=0$ and $X\ne0$.
Alignment makes $dp$ parallel to $X$, and
\eqref{stat:global-flux-laws} gives $\mathcal D\cdot dp=0$.
Thus~\eqref{stat:polynomial-energy} holds there as well, with
equality in its second assertion.

The current $j$ is bounded up to $S$.  This is immediate when
$u>0$ there.  In the other case $dW=O(s)$ and
$X^j(dX^\flat)_{ij}=O(s^3)$, so the quotient
in~\eqref{stat:global-Killing-flux} extends smoothly by
$u=sa$.  The other current is
$u^{-1}\nabla p=(k-1)E$, already smooth.  Since $P|_S=0$, a
distance cutoff gives a boundary error tending to zero.  At
infinity,
\[
 P=O_2(r^{-q_0}),\qquad j=O(r^{-1-q_0}),
\]
so the cutoff error in integrating $P\operatorname{div}j$ is
$O(R^{n-2-2q_0})$.  Testing with the product of the boundary and
end cutoffs used above, and using nonnegativity in
\eqref{stat:polynomial-energy}, gives
\[
 \int_{\Omega^\circ}u^{-1}|dP|^2=0.
\]
Thus $P$ is constant, and its limit at infinity is zero.  This proves
the first equality in~\eqref{stat:W-polynomial}.

Its root at $p_b$ is simple.  Otherwise $W=(p-p_b)^2$.  If $u>0$
on $S$, this contradicts $dW=2\kappa X^\flat\ne0$ there.  If
$u=0$ on $S$, then $dp=(k-1)uE^\flat$ gives
$p-p_b=O(s^2)$, whereas~\eqref{stat:zero-lapse-factors} gives
$W=\kappa^2s^2+O(s^3)$, again a contradiction.
Writing $p=t p_b$ with $0<t<1$, the polynomial becomes
\[
 W=(1-t)(1-p_b^2t).
\]
Every intermediate value is attained by continuity between the
boundary and the end.  Nonnegativity implies $p_b^2\leq1$, and
simplicity excludes equality.  The displayed formula then makes
$W$ strictly positive at every interior point.
\end{proof}

\begin{lemma}[Regular attachment of the orbit base]
\label{stat:base-attachment}
On each positive component adjacent to $S$, the pair $(h_s,\lambda)$
has a smooth nondegenerate attachment at $S$.  The attached boundary
metric is $g|_{TS}$, and
\begin{equation}
 \lambda|_S=0,\qquad |d\lambda|_{h_s}|_S=\kappa.
 \label{stat:base-boundary-gradient}
\end{equation}
If $u>0$ on $S$, the attachment uses coordinates $(\lambda,y)$ in
place of original coordinates $(W,y)$, and
\begin{equation}
 X^\flat=\frac1{2\kappa}dW+W\beta
 \label{stat:positive-boundary-one-form}
\end{equation}
for a smooth one-form $\beta$ in the original collar.  In these
attached coordinates $h_s$ is smooth and invariant, as a tensor,
under $\lambda\mapsto-\lambda$.  If $u=0$ on $S$, the original
smooth structure is retained, $\lambda$ is a smooth defining
function, and $A_s$ is smooth up to $S$.
\end{lemma}

\begin{proof}
In the positive-$u$ case, $W$ is an original boundary defining
function.  Equation~\eqref{stat:positive-boundary-one-form} follows
from the boundary identity $dW=2\kappa X^\flat$ and smooth
division by this defining function.  More explicitly, write
\[
 X^\flat=b(W,y)dW+W\beta_A(W,y)dy^A,
 \qquad b(0,y)=\frac1{2\kappa}.
\]
Upon setting $W=\lambda^2$, the metric components are
\begin{align*}
 (h_s)_{\lambda\lambda}
   &=4\lambda^2g_{WW}+4b^2,\\
 (h_s)_{\lambda A}
   &=2\lambda(g_{WA}+b\beta_A),\\
 (h_s)_{AB}&=g_{AB}+\lambda^2\beta_A\beta_B,
\end{align*}
where every coefficient on the right is evaluated at
$(\lambda^2,y)$.  Normal and tangential diagonal coefficients are
even and the mixed coefficients are odd, which is exactly tensor
invariance under reflection.  At $\lambda=0$ the normal coefficient
is $\kappa^{-2}$, the mixed coefficients vanish, and the tangential
metric is $g|_{TS}$.

In the zero-$u$ case, put
$v=(a^2-s^2|Y|^2)^{1/2}$ in a small original collar.  Then
\[
 \lambda=sv,\qquad
 h_s=g+s^2v^{-2}Y^\flat\otimes Y^\flat,\qquad
 A_s=v^{-2}Y^\flat.
\]
These expressions are smooth and nondegenerate, and $v|_S=\kappa$.
They prove the remaining assertions.
\end{proof}

\begin{lemma}[Local static electrovacuum equations]
\label{stat:static-equations}
On each positive component, the extra matter vanishes and the
stationary splitting is locally static.  Its field is
\begin{equation}
 F=\frac{c_n}{k-1}(dz-A_s)\wedge dp.
 \label{stat:static-F}
\end{equation}
The orbit fields satisfy
\begin{align}
 \Delta_{h_s}\lambda&=\lambda^{-1}|dp|_{h_s}^2,
 &\operatorname{div}_{h_s}(\lambda^{-1}\nabla_{h_s}p)&=0,
 \label{stat:static-lapse-Maxwell}\\
 \operatorname{Ric}_{h_s}-\lambda^{-1}\operatorname{Hess}_{h_s}\lambda
   &=\frac{-k\,dp\otimes dp+|dp|_{h_s}^2h_s}
           {(k-1)\lambda^2}.
 \label{stat:static-Ricci}
\end{align}
Equivalently, with $\check h=\lambda^{2/(k-1)}h_s$,
\begin{align}
 \operatorname{Ric}_{\check h}
   &=\frac{k}{k-1}
       \left(d\log\lambda\otimes d\log\lambda
                       -\lambda^{-2}dp\otimes dp\right),
 \label{stat:conformal-static-Ricci}\\
 \operatorname{div}_{\check h}
       (\lambda^{-2}\nabla_{\check h}p)&=0.
 \label{stat:conformal-Maxwell}
\end{align}
\end{lemma}

\begin{proof}
Where $W>0$, equation~\eqref{stat:remaining-matter} has no matter
term.  Lemma~\ref{stat:vanishing-twist} makes $A_s$ locally exact,
so a local coordinate $T$ has $dT=dz-A_s$ and
$\mathbf g=-\lambda^2dT^2+h_s$.  Alignment gives
$A_s\wedge dp=0$, proving~\eqref{stat:static-F}.
In this splitting the normalized static electric covector is
$dp/((k-1)\lambda)$.

The warped-product Ricci components are
\[
\begin{split}
 \operatorname{Ric}_{\mathbf g}(\partial_T,\partial_T)
 &=\lambda\Delta_{h_s}\lambda,\\
 \operatorname{Ric}_{\mathbf g,ij}
 &=\operatorname{Ric}_{h_s,ij}
 -\lambda^{-1}\operatorname{Hess}_{h_s,ij}\lambda.
\end{split}
\]
For $F_{Ti}=c_n p_i/(k-1)$, trace reversal gives respectively
$|dp|_{h_s}^2$ and the right side of
\eqref{stat:static-Ricci}.  The spatial divergence of the raised
$iT$ Maxwell component gives the second equation of
\eqref{stat:static-lapse-Maxwell}.  These calculations establish
all constants in~\eqref{stat:static-lapse-Maxwell}--
\eqref{stat:static-Ricci}.  Finally apply the conformal Ricci and
Laplacian formulas with conformal logarithm
$(k-1)^{-1}\log\lambda$.  The Hessian terms cancel and yield
\eqref{stat:conformal-static-Ricci}--
\eqref{stat:conformal-Maxwell}.
\end{proof}

\begin{proposition}[Charged reduction to vacuum]
\label{stat:charged-vacuum}
Suppose $Q\ne0$.  On the complete attached orbit base define
\begin{equation}
 \lambda_0=\frac{\lambda}{1-p_bp},\qquad
 h_0=(1-p_bp)^{2/(k-1)}h_s.
 \label{stat:explicit-vacuum-transform}
\end{equation}
Then $0<\lambda_0<1$ in the interior and
\begin{equation}
 \operatorname{Ric}_{h_0}
      =\lambda_0^{-1}\operatorname{Hess}_{h_0}\lambda_0,
 \qquad \Delta_{h_0}\lambda_0=0.
 \label{stat:vacuum-equations}
\end{equation}
The metric and lapse extend smoothly to the compact boundary, where
\begin{equation}
 \lambda_0=0,\qquad
 |d\lambda_0|_{h_0}
   =\kappa_0:=\kappa(1-p_b^2)^{-k/(k-1)}>0.
 \label{stat:vacuum-boundary-gradient}
\end{equation}
The boundary is totally geodesic and
$\operatorname{Hess}_{h_0}\lambda_0=0$ there.  The base is complete
including its boundary and has exactly the original coordinate end.
\end{proposition}

\begin{proof}
Lemma~\ref{stat:charged-polynomial} already gives $W>0$ throughout
the interior.  Set
\[
 a_0=p_b+d_0/2=\frac{p_b^2-1}{2p_b}\ne0,
 \qquad
 V_0=\frac12\log
      \frac{1-p/p_b}{1-p_bp}.
\]
The logarithm is real and $V_0\to0$ at infinity.  Direct
differentiation, using~\eqref{stat:W-polynomial}, gives
\begin{equation}
 dV_0=a_0\frac{dp}{W},\qquad
 d\log\lambda\otimes d\log\lambda-W^{-1}dp\otimes dp
       =dV_0\otimes dV_0.
 \label{stat:vacuum-target-identity}
\end{equation}
The second identity follows also from
$(p+d_0/2)^2-W=d_0^2/4-1=a_0^2$.
Equations~\eqref{stat:conformal-static-Ricci} and
\eqref{stat:conformal-Maxwell} now give
\[
 \operatorname{Ric}_{\check h}
      =\frac{k}{k-1}dV_0\otimes dV_0,
 \qquad \Delta_{\check h}V_0=0.
\]
The inverse conformal formulas, with
$\lambda_0=e^{V_0}$ and
$h_0=\lambda_0^{-2/(k-1)}\check h$, prove
\eqref{stat:vacuum-equations}.  The polynomial factorization makes
these definitions exactly~\eqref{stat:explicit-vacuum-transform}.
It also gives $0<\lambda_0<1$, since for $p=t p_b$, $0<t<1$,
\[
 \lambda_0^2=\frac{1-t}{1-p_b^2t}<1.
\]

Near the boundary the simple root in
Lemma~\ref{stat:charged-polynomial} makes $p$ a smooth function of
$W=\lambda^2$.  Thus the positive factor $1-p_bp$ is smooth in the
attached coordinates of Lemma~\ref{stat:base-attachment}.  Its
boundary value is $1-p_b^2>0$, and the normal rescaling gives
\eqref{stat:vacuum-boundary-gradient}.  Smoothness of $h_0$ and
the equation
$\operatorname{Hess}_{h_0}\lambda_0
=\lambda_0\operatorname{Ric}_{h_0}$ imply that the Hessian
vanishes at the boundary.  Its tangential components are the
nonzero normal derivative times the boundary second fundamental
form, so the boundary is totally geodesic.

Completeness of $h_s$ including its attachment follows from
$h_s\geq g$.  Indeed a finite-length Cauchy escape has a limit in
the complete original metric space.  An interior limit is harmless
because $W>0$ there, and a boundary limit is contained in the
regular attached collar.  The same reasoning proves metric
completeness, or equivalently the absence of any finite-length
escape.  Since
\[
 1-p_b^2\leq1-p_bp\leq1,
\]
$h_0$ and $h_s$ are uniformly comparable.  Completeness passes to
$h_0$.  The attachment changes only the boundary collar's smooth
defining function, so the topological end and compact-complement
properties are retained.
\end{proof}

\subsection{Zero charge and exclusion of an interior frontier}

\begin{proposition}[Global positivity when the charge vanishes]
\label{stat:zero-charge-positivity}
If $Q=0$, then $E=0$, $p=0$, and $W>0$ on all of
$\Omega^\circ$.  The attached pair $(h_s,\lambda)$ is a complete
vacuum static pair with $0<\lambda<1$, connected compact boundary,
and boundary derivative $\kappa$.
\end{proposition}

\begin{proof}
The boundary value of $p$ is constant by
Lemma~\ref{stat:electric-alignment}.  Integrate
$\operatorname{div}(pE)=dp(E)=(k-1)u|E|^2$ on the original
exterior truncated at radius $R$.  The boundary normal in this
divergence theorem is $-\nu$ on $S$.  The outer term tends to zero
by the decay of $p$ and $E$, and the inner flux is
$-p_b\omega Q=0$.  Since $u$ is bounded and $|E|^2$ is integrable,
this gives
\[
 (k-1)\int_\Omega u|E|^2=0.
\]
Positivity of $u$ in the interior implies $E=0$, and then $p=0$ by
its normalization at infinity.

On every component of $\mathcal U$, Lemmas
\ref{stat:vanishing-twist} and~\ref{stat:static-equations} now give
\begin{equation}
 \operatorname{Hess}_{h_s}\lambda
     =\lambda\operatorname{Ric}_{h_s},\qquad
 \Delta_{h_s}\lambda=0,
 \qquad \operatorname{Scal}_{h_s}=0.
 \label{stat:local-vacuum}
\end{equation}
Let $\mathcal C$ be the positive component containing the whole
sufficiently distant end.  Any other positive component has compact
closure in the original exterior and continuous boundary value
$\lambda=0$ on its whole topological boundary, including any part
on $S$.  Its positive maximum would be attained in its interior,
contradicting the strong maximum principle.  Thus no other positive
component exists, and all deleted boundary collars lie in
$\mathcal C$.  The same compact-superlevel argument gives
$0<\lambda<1$ on $\mathcal C$.

We next exclude a finite orbit-metric distance to an interior
zero.  Near any such zero, $u$ and $|X|$ are bounded below, and
Lemma~\ref{stat:vanishing-twist} makes the left side of
\eqref{stat:exact-twist-flux} zero.  Therefore
\begin{equation}
 (d_T\log W)_i=|X|^{-2}X^j(dX^\flat)_{ij}.
 \label{stat:zero-charge-transverse-log}
\end{equation}
Its right side is bounded.  Together with~\eqref{stat:glaeser},
this gives on $\mathcal C$ near the zero
\begin{equation}
 |d\log W|_{h_s}^2
 =|d_T\log W|_g^2
      +\frac{(dW(e))^2}{u^2W}\leq C.
 \label{stat:zero-charge-distance}
\end{equation}
Consequently a finite $h_s$-length curve cannot approach the zero:
integration of~\eqref{stat:zero-charge-distance} would keep
$\log W$ bounded on its tail.

Attach $S$ to the base of $\mathcal C$ as in
Lemma~\ref{stat:base-attachment}.  By~\eqref{stat:local-vacuum},
the Hessian of $\lambda$ vanishes at $S$, and the boundary is
totally geodesic.  In Gaussian base distance $\sigma$ from $S$,
smoothness gives the Taylor expansions
\begin{align}
 h_s&=d\sigma^2+
       \bigl(\gamma_0+\sigma^2\gamma_2
                         +\sigma^3R\bigr)_{AB}dy^A dy^B,
 \label{stat:bolt-base-Taylor}\\
 \lambda&=\kappa\sigma+a_3(y)\sigma^3
                              +\sigma^4R_\lambda.
 \label{stat:bolt-lapse-Taylor}
\end{align}
The remainders are smooth for $\sigma\geq0$.  The absent quadratic
term in the lapse is exactly
$\operatorname{Hess}_{h_s}\lambda(\partial_\sigma,
\partial_\sigma)|_S=0$.

Take a circle coordinate $T$ of period $2\pi/\kappa$, and put
\begin{equation}
 h_{\mathrm{circ}}=h_s+\lambda^2dT^2
 \quad\hbox{on }\mathcal C\times
                  (\mathbb R/(2\pi/\kappa)\mathbb Z).
 \label{stat:circle-metric}
\end{equation}
Collapse this circle at $S$ by attaching a disk in the normal
direction.  The period is the same over the entire connected
boundary because $\kappa$ is constant.  To check regularity, set
$\theta=\kappa T$, use disk coordinates
$(x_1,x_2)=(\sigma\cos\theta,\sigma\sin\theta)$, and write
$\alpha_D=x_1dx_2-x_2dx_1$.  Besides the Euclidean disk metric,
the polar contribution is
\begin{align*}
 \left[\left(\frac{\lambda}{\kappa\sigma}\right)^2-1\right]
        \sigma^2d\theta^2
 &=\left[\left(\frac{\lambda}{\kappa\sigma}\right)^2-1\right]
        \frac{\alpha_D^2}{\sigma^2}\\
 &=\frac{2a_3(y)}\kappa\alpha_D^2
          +\sigma B(\sigma,y)\alpha_D^2
\end{align*}
with smooth $B$.  The first term has smooth quadratic coefficients
in $(x_1,x_2)$; the remaining coefficients have differentiated
order at least three.  The same holds for the remainder in
\eqref{stat:bolt-base-Taylor}.  The extended metric is therefore
$C^{2,\alpha}$ for every $0<\alpha<1$.  It is nondegenerate.
The warped-product Ricci formulas and~\eqref{stat:local-vacuum}
give Ricci curvature zero off the attached copy of $S$, and
continuity extends this identity across it.  Harmonic-coordinate
regularity for Einstein metrics then makes this extension smooth
\cite{DeTurckKazdan1981}.

The resulting manifold, denoted by $\mathcal M$, is complete and
has no boundary.  To check completeness directly, a finite-length
curve projects to a finite $h_s$-length curve, hence to a
$g$-Cauchy curve because $h_s\geq g$.  Completeness of the original
metric space supplies a limit in $\Omega$.  An interior zero of
$W$ is excluded by~\eqref{stat:zero-charge-distance}.  At an
interior point of $\mathcal C$, the circle has uniformly positive
size, so its coordinate is Cauchy as well.  At a point of $S$, the
regular disk attachment makes all limiting circle phases converge
to that point of the attached zero section.  Thus no finite-length
escape remains, which proves completeness.

Suppose an interior frontier of $\mathcal C$ still existed.  A
sequence in $\mathcal C$ approaching such a frontier, lifted at a
fixed circle phase, escapes every compact subset of $\mathcal M$.
Choose a far coordinate sphere and its product with the circle,
\[
 \Sigma_R\cong S^{n-1}\times S^1.
\]
It is a compact connected cross-section.  Its outside is the
asymptotic product end, while its inside contains the escaping
sequence and is noncompact.  The inside is connected: excursions
of a path through the outside can be replaced by paths in a thin
inside collar of the connected cross-section.  A compact tubular
neighborhood of $\Sigma_R$ therefore separates two noncompact
regions.  In particular $\mathcal M$ has at least two ends.

For clarity, the line needed for the splitting theorem can be
constructed directly.  Choose escaping sequences in the two
regions and minimizing segments between them.  Each segment meets
the fixed compact separating cross-section, and its distances from
that cross-section to both endpoints tend to infinity.  Centering
there and taking a subsequence gives a complete minimizing line.
The smooth, complete, connected manifold $\mathcal M$ has
$\operatorname{Ric}=0$, so the Cheeger--Gromoll splitting theorem
applies \cite{CheegerGromoll1971}.  It gives
$\mathcal M\cong\mathbb R\times N$ with $N$ complete and
connected.  Since the product with a connected noncompact factor
has only one end, $N$ must be compact.  Such a product has at most
linear volume growth.

On the other hand,~\eqref{stat:end-interface} shows that $h_s$
approaches a positive constant Euclidean metric and $\lambda\to1$.
The circle has fixed period.  Hence the product end in
\eqref{stat:circle-metric} gives
$\operatorname{Vol}_{h_{\mathrm{circ}}}(B_R)\geq cR^n$ for large
$R$: paths to a fixed far sphere have bounded cost, and the metric
on the remaining end is uniformly comparable to its Euclidean
product model.  This contradicts linear growth.  There is no
interior frontier.  Since $\Omega^\circ$ is connected,
$\mathcal C=\Omega^\circ$ and $W>0$ throughout.

Now $h_s$ itself is complete up to its regular boundary by the
same projection and collar argument, and
\eqref{stat:local-vacuum} holds globally.  The boundary assertions
follow from Lemma~\ref{stat:base-attachment} and the vanishing
Hessian already proved.
\end{proof}

\begin{proposition}[Complete vacuum reduction]
\label{stat:vacuum-reduction}
Under the stationary interface stated at the start of this section,
the Killing field is timelike on the whole open exterior:
$W>0$ on $\Omega^\circ$.  The original matter densities vanish
there, $X\wedge E=0$, and $dA_s=0$.  The orbit pair
$(h_s,\lambda)$ has the regular attachment in
Lemma~\ref{stat:base-attachment}.

If $Q\ne0$, define $(h_0,\lambda_0)$ by
\eqref{stat:explicit-vacuum-transform}; if $Q=0$, take
$(h_0,\lambda_0)=(h_s,\lambda)$.  In either case this is a smooth
vacuum static pair satisfying~\eqref{stat:vacuum-equations} on a
connected orientable exterior, complete including its connected
compact boundary, with compact complement of exactly one Euclidean
coordinate end.  Its lapse is zero at the boundary, lies strictly
between zero and one in the interior, and tends to one at infinity.
The inward boundary derivative is a positive constant, equal to
\eqref{stat:vacuum-boundary-gradient} in the charged case and to
$\kappa$ in the uncharged case.  The boundary is totally geodesic
and $\operatorname{Hess}_{h_0}\lambda_0=0$ there.  Gaussian
reflection of $h_0$ with odd reflection of $\lambda_0$ is at least
$C^{2,\alpha}$ for $0<\alpha<1$.

After a fixed linear change of end coordinates,
\begin{equation}
 h_0-\delta=O_2(r^{-q_0}),\qquad
 \lambda_0-1=O_2(r^{-q_0})
 \label{stat:vacuum-end}
\end{equation}
for $q_0$ as in~\eqref{stat:end-exponent}.  No simple-connectivity
claim or global primitive of $A_s$ is needed for this reduction.
\end{proposition}

\begin{proof}
All interior, boundary, and completeness assertions follow from
Propositions~\ref{stat:charged-vacuum} and
\ref{stat:zero-charge-positivity}, together with
Lemmas~\ref{stat:electric-alignment} and
\ref{stat:vanishing-twist}.  On $W>0$,
\eqref{stat:remaining-matter} gives $\mu_m=0$, and the matter
dominant energy condition gives $J_m=0$.  Continuity gives the
same conclusion up to $S$.

In the original end coordinates, $h_s$ tends to
$\delta+b^\flat\otimes b^\flat$, a positive constant matrix.
Equations~\eqref{stat:end-interface} and
\eqref{stat:orbit-quantities} give differentiated error
$O_2(r^{-q_0})$.  In the charged case the additional conformal
factor differs from one by $O_2(r^{2-n})$, so the same bound holds
for $h_0$ and $\lambda_0$.  A fixed linear coordinate change
normalizes the limiting matrix and preserves these decay orders.

Finally take Gaussian normal coordinates for $h_0$ at the boundary.
Total geodesicity makes the first normal derivative of its
tangential metric zero.  The vanishing boundary Hessian makes the
second normal derivative of the lapse zero.  Even metric reflection
and odd lapse reflection consequently match their derivatives
through order two and are $C^{2,\alpha}$.  This is the regularity
required for the vacuum doubling argument in the next section.
\end{proof}
